\subsection{National Policy Guidelines for AI Development, Employment, and Antifascism}
\label{sec:8.4.1}
A national AI policy has a predictable shape: fund value-alignment research and the work of operationalizing ethical principles inside a system; build standing interdisciplinary capacity; train a workforce; consult the public through assemblies, consultations, and partnerships with organizations outside government; pay for robustness and long-term safety research; and update data-privacy, cybersecurity, accountability and legal-status law where AI has outrun it, on which sections~\ref{sec:8.4.4} and \ref{sec:8.6} carry the specifics. Every country that has written a national strategy has written most of that list. The list is not where national policies differ.

They differ in two places, both having to do with leverage.

The first is harmonization. A regulation with teeth in one country does little if a developer can relocate to one without them, which makes a national framework partly an instrument for binding other national frameworks — through mutual recognition, through market access, through the cost of maintaining separate compliance regimes side by side. Sections~\ref{sec:8.7.4} and \ref{sec:8.7.8} work through how little of that has actually happened and which mechanisms have carried anything at all across a border. A policy written as though the country were alone is a policy about domestic firms, and the firms that matter most are not domestic anywhere.

The second is what a government adds to industry self-regulation. Section~\ref{sec:8.4.2} spends its length on why a voluntary code is worth less than it announces. What national policy contributes is exactly the part self-regulation lacks by construction: a standard or certification program backed by government monitoring carries a cost for noncompliance, where a purely voluntary one carries a cost only to reputation, which the company itself gets to price.

Neither of these is expensive. Neither needs a new agency, a new science, or an argument the public has not already accepted. That both remain the exception is a fact about political attention, and not about difficulty.
